\newif\ifglobal

%%%%%%%%%%%%%%%%%%%%%%%
%%% Compile to view %%%
%%%%%%%%%%%%%%%%%%%%%%%

%%%%%%%%%%%%%%%%%%%%%%%%%%%%%%%%%%%%%%%%%%%%%%%%%%%
%%% Readme:                                     %%%
%%% https://www.overleaf.com/read/bwdjvfjksvjy/ %%%
%%%%%%%%%%%%%%%%%%%%%%%%%%%%%%%%%%%%%%%%%%%%%%%%%%%

% >>> M?: marks a module setting number or important notice

% >>> M4: Set {./relative-path} to external files for this module <<<
\def \modulefiles {./30-AES}
% To add an external file, use:
% (Replace '---' with actual filename)
% '\input{\modulefiles/tables/---.tex}' for tables
% '\includegraphics{\modulefiles/figures/---}' for figures
% '\subfile{\modulefiles/---}' for register files

% >>> M1: This module is in global project? <<<
% 'YES' - uncomment; 'NO' - comment out
\globaltrue

\ifglobal
    \documentclass[main\_\_EN.tex]{subfiles}

\else
    \documentclass[openany, 10pt]{book}

    % >>> M2: In ./00-shared/preamble-trm-module, also find
    %         settings for visibility of todo notes, watermark <<<
    \usepackage{preamble-trm-module}

    % >>> M3: Temporary labels in English,
    %         you can add your own labels there <<<
    \myexternaldocument{./00-shared/temp-labels-en}

\fi %\ifglobal


\setcounter{secnumdepth}{4}

% >>> M5: If this module needs any updates in
% './00-shared/preamble-trm-module.sty'
% (include additional package, etc.), contact your module PM
% or create the file './00-shared/changelog.tex' make the
% necessary updates yourself and document them in this file <<<


\begin{document}


% Sets language into which pre-defined bits of text are translated
\selectlanguage{english}

% Do not delete this block
\ifglobal
% Add the following front matter if
% this module is outside of global project
\else
    \InputIfFileExists{readme}{}{}
    {\let\clearpage\relax \listoftodos}
    \tableofcontents
    %\listoftables
    %\listoffigures
\fi


%%%%%%%%%%%%%%%%%%%%%%%%%
%%% TRM Module Begins %%%
%%%%%%%%%%%%%%%%%%%%%%%%%

\hypertarget{aes}{}
\chapter{AES Accelerator (AES)}\label{mod:aes}

\section{Introduction}

\chipname{} integrates an Advanced Encryption Standard (AES) Accelerator, which is a hardware device that speeds up AES Algorithm significantly, compared to AES algorithms implemented solely in software. The AES Accelerator integrated in \chipname{} has two working modes, which are \hyperref[sec:aes-typical-modes]{Typical AES} and \hyperref[sec:aes-dma-aes]{DMA-AES}.

\section{Features}

The following functionality is supported:
\begin{itemize}
    \item Typical AES working mode
        \begin{itemize}
            \item AES-128/AES-256 encryption and decryption
        \end{itemize}
    \end{itemize}

    \begin{itemize}
    \item DMA-AES working mode
        \begin{itemize}
            \item AES-128/AES-256 encryption and decryption
            \item Block cipher mode
                \begin{itemize}
                    \item ECB (Electronic Codebook)
                    \item CBC (Cipher Block Chaining)
                    \item OFB (Output Feedback)
                    \item CTR (Counter)
                    \item CFB8 (8-bit Cipher Feedback)
                    \item CFB128 (128-bit Cipher Feedback)
                \end{itemize}
            \item Interrupt on completion of computation
        \end{itemize}
    \end{itemize}


\section{AES Working Modes}

The AES Accelerator integrated in \chipname{} has two working modes, which are \hyperref[sec:aes-typical-modes]{Typical AES} and \hyperref[sec:aes-dma-aes]{DMA-AES}.

\begin{itemize}
    \item Typical AES Working Mode: \begin{itemize}
        \item Supports encryption and decryption using cryptographic keys of 128 and 256 bits, specified in \href{https://nvlpubs.nist.gov/nistpubs/FIPS/NIST.FIPS.197.pdf}{NIST FIPS 197}.
    \end{itemize}
    In this working mode, the plaintext and ciphertext is written and read via CPU directly.
    \item DMA-AES Working Mode:
    \begin{itemize}
        \item Supports encryption and decryption using cryptographic keys of 128 and 256 bits, specified in \href{https://nvlpubs.nist.gov/nistpubs/FIPS/NIST.FIPS.197.pdf}{NIST FIPS 197};
        \item Supports block cipher modes ECB/CBC/OFB/CTR/CFB8/CFB128 under \href{https://nvlpubs.nist.gov/nistpubs/Legacy/SP/nistspecialpublication800-38a.pdf}{NIST SP 800-38A}.
    \end{itemize}
    In this working mode, the plaintext and ciphertext is written and read via DMA. An interrupt will be generated when operation completes.
\end{itemize}

Users can choose the working mode for AES accelerator by configuring the \hyperref[regdesc:AESDMAENABLEREG]{AES\_DMA\_ENABLE\_REG} register according to Table \ref{tab:aes-dma-enable} below.

\begin{longtable}[c]{ | c | l | }
\caption{AES Accelerator Working Mode}
    \label{tab:aes-dma-enable}
    \\\hline
    \rowcolor{lightgray}
\hyperref[regdesc:AESDMAENABLEREG]{AES\_DMA\_ENABLE\_REG}       & Working Mode\\
        \hline
        0                       & Typical AES    \\
        \hline
        1                       & DMA-AES        \\
        \hline
\end{longtable}

Users can choose the length of cryptographic keys and encryption/decryption by configuring the \hyperref[regdesc:AESMODEREG]{AES\_MODE\_REG} register according to Table \ref{tab:aes-operation} below.

\begin{longtable}[c]{ | c | l | }
\caption{Key Length and Encryption/Decryption}
\label{tab:aes-operation}
\\\hline
\rowcolor{lightgray}
\hyperref[regdesc:AESMODEREG]{AES\_MODE\_REG}[2:0]     & Key Length and Encryption/Decryption\\
    \hline
    0                   & AES-128 encryption    \\
    \hline
    1                   & reserved    \\
    \hline
    2                   & AES-256 encryption    \\
    \hline
    3                   & reserved    \\
    \hline
    4                   & AES-128 decryption    \\
    \hline
    5                   & reserved    \\
    \hline
    6                   & AES-256 decryption    \\
    \hline
    7                   & reserved    \\
    \hline
\end{longtable}

For detailed introduction on these two working modes, please refer to Section \ref{sec:aes-typical-modes} and Section \ref{sec:aes-dma-aes} below.

\begin{importantlisting}
\vspace{-0.5em}
    \chipname{}'s \nameref{mod:digsig} module will call the AES accelerator. Therefore, users cannot access the AES accelerator when \nameref{mod:digsig} module is working.
\end{importantlisting}

\section{Typical AES Working Mode}\label{sec:aes-typical-modes}

In the Typical AES working mode,
users can check the working status of the AES accelerator by inquiring the \hyperref[regdesc:AESSTATEREG]{AES\_STATE\_REG} register and comparing the return value against the Table \ref{tab:aes-state-value-typical} below.

\begin{longtable}[c]{ | c | l | l| }
\caption{Working Status under Typical AES Working Mode}
\label{tab:aes-state-value-typical}
\\\hline
\rowcolor{lightgray}
\hyperref[regdesc:AESSTATEREG]{AES\_STATE\_REG} & Status &  Description \\ \hline

    0                   & IDLE      &   The AES accelerator is idle or completed operation.\\ \hline
    1                   & WORK      &   The AES accelerator is in the middle of an operation.     \\ \hline
\end{longtable}

\subsection{Key, Plaintext, and Ciphertext}

The encryption or decryption key is stored in \hyperref[regdesc:AESKEYnREG]{AES\_KEY\_\regindex{n}\_REG}, which is a set of eight 32-bit registers.

\begin{itemize}
    \item For AES-128 encryption/decryption, the 128-bit key is stored in \hyperref[regdesc:AESKEYnREG]{AES\_KEY\_0\_REG} \verb+~+ \hyperref[regdesc:AESKEYnREG]{AES\_KEY\_3\_REG}.
    \item For AES-256 encryption/decryption, the 256-bit key is stored in \hyperref[regdesc:AESKEYnREG]{AES\_KEY\_0\_REG} \verb+~+ \hyperref[regdesc:AESKEYnREG]{AES\_KEY\_7\_REG}.
\end{itemize}

The plaintext and ciphertext are stored in \hyperref[regdesc:AESTEXTINnREG]{AES\_TEXT\_IN\_\regindex{m}\_REG} and \hyperref[regdesc:AESTEXTOUTnREG]{AES\_TEXT\_OUT\_\regindex{m}\_REG}, which are two sets of four 32-bit registers.

\begin{itemize}
    \item For AES-128/AES-256 encryption, the \hyperref[regdesc:AESTEXTINnREG]{AES\_TEXT\_IN\_\regindex{m}\_REG} registers are initialized with plaintext. Then, the AES Accelerator stores the ciphertext into \hyperref[regdesc:AESTEXTOUTnREG]{AES\_TEXT\_OUT\_\regindex{m}\_REG} after operation.
    \item For AES-128/AES-256 decryption, the \hyperref[regdesc:AESTEXTINnREG]{AES\_TEXT\_IN\_\regindex{m}\_REG} registers are initialized with ciphertext. Then, the AES Accelerator stores the plaintext into \hyperref[regdesc:AESTEXTOUTnREG]{AES\_TEXT\_OUT\_\regindex{m}\_REG} after operation.
\end{itemize}


\subsection{Endianness}

\textbf{Text Endianness}

In Typical AES working mode, the AES Accelerator uses cryptographic keys to encrypt and decrypt data in blocks of 128 bits.
When filling data into \hyperref[regdesc:AESTEXTINnREG]{AES\_TEXT\_IN\_\regindex{m}\_REG} register or reading result from \hyperref[regdesc:AESTEXTINnREG]{AES\_TEXT\_OUT\_\regindex{m}\_REG} registers, users should follow the text endianness type specified in Table \ref{tab:aes-endianness-type-typical-plaintext-ciphertext}.


\begin{table}[H]
    \centering
    \caption{Text Endianness Type for Typical AES}
    \label{tab:aes-endianness-type-typical-plaintext-ciphertext}

    \begin{threeparttable}

    \begin{tabular}{ |c|c|c|c|c|c|}
    \hline
        \rowcolor{lightgray}
        \multicolumn{6}{|c|}{Plaintext/Ciphertext} \\\hline
        \multicolumn{2}{|c|}{\multirow{2}{*}{State\tnote{1}}} &\multicolumn{4}{c|}{c\tnote{2}}\\ \cline{3-6}
        \multicolumn{2}{|c|}{} &0 &1 &2 &3\\ \cline{1-6}
        \multirow{4} {*}r & 0 &\small{\hyperref[regdesc:AESTEXTINnREG]{AES\_TEXT\_\regindex{x}\_0\_REG}[7:0]} &   \small{\hyperref[regdesc:AESTEXTINnREG]{AES\_TEXT\_\regindex{x}\_1\_REG}[7:0]}  &   \small{\hyperref[regdesc:AESTEXTINnREG]{AES\_TEXT\_\regindex{x}\_2\_REG}[7:0]}  &   \small{\hyperref[regdesc:AESTEXTINnREG]{AES\_TEXT\_\regindex{x}\_3\_REG}[7:0]}\\
        \cline{2-6}
        &1       &   \small{\hyperref[regdesc:AESTEXTINnREG]{AES\_TEXT\_\regindex{x}\_0\_REG}[15:8]} &   \small{\hyperref[regdesc:AESTEXTINnREG]{AES\_TEXT\_\regindex{x}\_1\_REG}[15:8]}     &   \small{\hyperref[regdesc:AESTEXTINnREG]{AES\_TEXT\_\regindex{x}\_2\_REG}[15:8]} &   \small{\hyperref[regdesc:AESTEXTINnREG]{AES\_TEXT\_\regindex{x}\_3\_REG}[15:8]}\\
        \cline{2-6}
        &2   &   \small{\hyperref[regdesc:AESTEXTINnREG]{AES\_TEXT\_\regindex{x}\_0\_REG}[23:16]}    &   \small{\hyperref[regdesc:AESTEXTINnREG]{AES\_TEXT\_\regindex{x}\_1\_REG}[23:16]}    &   \small{\hyperref[regdesc:AESTEXTINnREG]{AES\_TEXT\_\regindex{x}\_2\_REG}[23:16]}    &   \small{\hyperref[regdesc:AESTEXTINnREG]{AES\_TEXT\_\regindex{x}\_3\_REG}[23:16]}\\
        \cline{2-6}
        &3   &   \small{\hyperref[regdesc:AESTEXTINnREG]{AES\_TEXT\_\regindex{x}\_0\_REG}[31:24]}    &   \small{\hyperref[regdesc:AESTEXTINnREG]{AES\_TEXT\_\regindex{x}\_1\_REG}[31:24]}    & \small{\hyperref[regdesc:AESTEXTINnREG]{AES\_TEXT\_\regindex{x}\_2\_REG}[31:24]}  &   \small{\hyperref[regdesc:AESTEXTINnREG]{AES\_TEXT\_\regindex{x}\_3\_REG}[31:24]}\\
        \hline
        \end{tabular}

        \begin{tablenotes}
            \item[1] The definition of ``State (including c and r)" is described in Section 3.4 The State in \href{https://nvlpubs.nist.gov/nistpubs/FIPS/NIST.FIPS.197.pdf}{NIST FIPS 197}.
            \item[2] Where \regindex{x} = IN or OUT.
        \end{tablenotes}
     \end{threeparttable}
    \end{table}


\textbf{Key Endianness}

In Typical AES working mode,
When filling key into \hyperref[regdesc:AESKEYnREG]{AES\_KEY\_\regindex{m}\_REG} registers, users should follow the key endianness type specified in Table \ref{tab:aes-128-key-endian} and Table \ref{tab:aes-256-key-endian}.


\begin{table}[H]
    \centering
    \caption{Key Endianness Type for AES-128 Encryption and Decryption}
    \label{tab:aes-128-key-endian}

    \begin{threeparttable}
    \begin{tabular}{ |l|l|l|l|l| }
    \hline
        \rowcolor{lightgray}
        Bit\tnote{1} &   w[0]    &   w[1]    &   w[2]    &   w[3]\tnote{2}     \\\hline

        % table contents
        [31:24]     &   \hyperref[regdesc:AESKEYnREG]{AES\_KEY\_0\_REG}[7:0]    &   \hyperref[regdesc:AESKEYnREG]{AES\_KEY\_1\_REG}[7:0]    &   \hyperref[regdesc:AESKEYnREG]{AES\_KEY\_2\_REG}[7:0]    &   \hyperref[regdesc:AESKEYnREG]{AES\_KEY\_3\_REG}[7:0]\\ \hline
        [23:16]     &   \hyperref[regdesc:AESKEYnREG]{AES\_KEY\_0\_REG}[15:8]   &   \hyperref[regdesc:AESKEYnREG]{AES\_KEY\_1\_REG}[15:8]   &   \hyperref[regdesc:AESKEYnREG]{AES\_KEY\_2\_REG}[15:8]   &   \hyperref[regdesc:AESKEYnREG]{AES\_KEY\_3\_REG}[15:8]\\ \hline
        [15:8]  &   \hyperref[regdesc:AESKEYnREG]{AES\_KEY\_0\_REG}[23:16]  &   \hyperref[regdesc:AESKEYnREG]{AES\_KEY\_1\_REG}[23:16]  &   \hyperref[regdesc:AESKEYnREG]{AES\_KEY\_2\_REG}[23:16]  &   \hyperref[regdesc:AESKEYnREG]{AES\_KEY\_3\_REG}[23:16]\\ \hline
        [7:0]   &   \hyperref[regdesc:AESKEYnREG]{AES\_KEY\_0\_REG}[31:24]  &   \hyperref[regdesc:AESKEYnREG]{AES\_KEY\_1\_REG}[31:24]  & \hyperref[regdesc:AESKEYnREG]{AES\_KEY\_2\_REG}[31:24]    &   \hyperref[regdesc:AESKEYnREG]{AES\_KEY\_3\_REG}[31:24]\\ \hline
        \end{tabular}

        \begin{tablenotes}
            \item[1] Column ``Bit" specifies the bytes of each word stored in w[0] \verb+~+ w[3].
            \item[2] w[0] \verb+~+ w[3] are ``the first Nk words of the expanded key" as specified in Section 5.2 Key Expansion in \href{https://nvlpubs.nist.gov/nistpubs/FIPS/NIST.FIPS.197.pdf}{NIST FIPS 197}.
        \end{tablenotes}
     \end{threeparttable}
\end{table}

\newgeometry{left=1cm,right=1cm,top=2cm,bottom=2cm}
\begin{landscape}


\begin{ThreePartTable}

\begin{longtable}[c]
{ |l|l|l|l|l|l|l|l|l|}
\caption{Key Endianness Type for AES-256 Encryption and Decryption}\label{tab:aes-256-key-endian}\\
\hline
\rowcolor{lightgray}
Bit$^1$ &   w[0]    &   w[1]    &   w[2]    &   w[3]    &   w[4]    &   w[5]&   w[6]    &   w[7]$^2$    \\\hline
[31:24]     &   \scriptsize \hyperref[regdesc:AESKEYnREG]{AES\_KEY\_0\_REG}[7:0]    &   \scriptsize \hyperref[regdesc:AESKEYnREG]{AES\_KEY\_1\_REG}[7:0]    &   \scriptsize \hyperref[regdesc:AESKEYnREG]{AES\_KEY\_2\_REG}[7:0]    &   \scriptsize \hyperref[regdesc:AESKEYnREG]{AES\_KEY\_3\_REG}[7:0]    &   \scriptsize \hyperref[regdesc:AESKEYnREG]{AES\_KEY\_4\_REG}[7:0]    &   \scriptsize \hyperref[regdesc:AESKEYnREG]{AES\_KEY\_5\_REG}[7:0]
& \scriptsize  \hyperref[regdesc:AESKEYnREG]{AES\_KEY\_6\_REG}[7:0]    &   \scriptsize \hyperref[regdesc:AESKEYnREG]{AES\_KEY\_7\_REG}[7:0]\\ \hline
[23:16]     &   \scriptsize \hyperref[regdesc:AESKEYnREG]{AES\_KEY\_0\_REG}[15:8]   &   \scriptsize \hyperref[regdesc:AESKEYnREG]{AES\_KEY\_1\_REG}[15:8]   &   \scriptsize \hyperref[regdesc:AESKEYnREG]{AES\_KEY\_2\_REG}[15:8]   &   \scriptsize \hyperref[regdesc:AESKEYnREG]{AES\_KEY\_3\_REG}[15:8]   &   \scriptsize \hyperref[regdesc:AESKEYnREG]{AES\_KEY\_4\_REG}[15:8]   &   \scriptsize \hyperref[regdesc:AESKEYnREG]{AES\_KEY\_5\_REG}[15:8]   &   \scriptsize \hyperref[regdesc:AESKEYnREG]{AES\_KEY\_6\_REG}[15:8]   &   \scriptsize \hyperref[regdesc:AESKEYnREG]{AES\_KEY\_7\_REG}[15:8]\\ \hline
[15:8]  &  \scriptsize \hyperref[regdesc:AESKEYnREG]{AES\_KEY\_0\_REG}[23:16]  &   \scriptsize \hyperref[regdesc:AESKEYnREG]{AES\_KEY\_1\_REG}[23:16]  &   \scriptsize \hyperref[regdesc:AESKEYnREG]{AES\_KEY\_2\_REG}[23:16]  &   \scriptsize \hyperref[regdesc:AESKEYnREG]{AES\_KEY\_3\_REG}[23:16]  &   \scriptsize \hyperref[regdesc:AESKEYnREG]{AES\_KEY\_4\_REG}[23:16]  &   \scriptsize \hyperref[regdesc:AESKEYnREG]{AES\_KEY\_5\_REG}[23:16]& \scriptsize \hyperref[regdesc:AESKEYnREG]{AES\_KEY\_6\_REG}[23:16]  &   \scriptsize \hyperref[regdesc:AESKEYnREG]{AES\_KEY\_7\_REG}[23:16]\\ \hline
[7:0]   &  \scriptsize \hyperref[regdesc:AESKEYnREG]{AES\_KEY\_0\_REG}[31:24]  &    \scriptsize \hyperref[regdesc:AESKEYnREG]{AES\_KEY\_1\_REG}[31:24]  & \scriptsize \hyperref[regdesc:AESKEYnREG]{AES\_KEY\_2\_REG}[31:24]    &   \scriptsize \hyperref[regdesc:AESKEYnREG]{AES\_KEY\_3\_REG}[31:24]  &   \scriptsize \hyperref[regdesc:AESKEYnREG]{AES\_KEY\_4\_REG}[31:24]  &   \scriptsize \hyperref[regdesc:AESKEYnREG]{AES\_KEY\_5\_REG}[31:24]& \scriptsize \hyperref[regdesc:AESKEYnREG]{AES\_KEY\_6\_REG}[31:24]  &   \scriptsize \hyperref[regdesc:AESKEYnREG]{AES\_KEY\_7\_REG}[31:24]\\ \hline
\end{longtable}
\begin{tablenotes}
\item[1] Column ``Bit" specifies the bytes of each word stored in w[0] \verb+~+ w[7].
\item[2] w[0] \verb+~+ w[7] are ``the first Nk words of the expanded key" as specified in Chapter 5.2 Key Expansion in \href{https://nvlpubs.nist.gov/nistpubs/FIPS/NIST.FIPS.197.pdf}{NIST FIPS 197}.
\end{tablenotes}
\end{ThreePartTable}

\end{landscape}
\restoregeometry


\subsection{Operation Process }

\textbf{Single Operation}

\begin{enumerate}
    \item Write 0 to the \hyperref[regdesc:AESDMAENABLEREG]{AES\_DMA\_ENABLE\_REG} register.
    \item Initialize registers \hyperref[regdesc:AESMODEREG]{AES\_MODE\_REG}, \hyperref[regdesc:AESKEYnREG]{AES\_KEY\_\regindex{n}\_REG}, \hyperref[regdesc:AESTEXTINnREG]{AES\_TEXT\_IN\_\regindex{m}\_REG}.
    \item Start operation by writing 1 to the \hyperref[regdesc:AESTRIGGERREG]{AES\_TRIGGER\_REG} register.
    \item Wait till the content of the \hyperref[regdesc:AESSTATEREG]{AES\_STATE\_REG} register becomes 0, which indicates the operation is completed.
    \item Read results from the \hyperref[regdesc:AESTEXTOUTnREG]{AES\_TEXT\_OUT\_\regindex{m}\_REG} register.
\end{enumerate}

\textbf{Consecutive Operations}

In consecutive operations, primarily the input \hyperref[regdesc:AESTEXTINnREG]{AES\_TEXT\_IN\_\regindex{m}\_REG} and output \hyperref[regdesc:AESTEXTOUTnREG]{AES\_TEXT\_OUT\_\regindex{m}\_REG} registers are being written and read, while the content of \hyperref[regdesc:AESDMAENABLEREG]{AES\_DMA\_ENABLE\_REG}, \hyperref[regdesc:AESMODEREG]{AES\_MODE\_REG}, \hyperref[regdesc:AESKEYnREG]{AES\_KEY\_\regindex{n}\_REG}
is kept unchanged. Therefore, the initialization can be simplified during the consecutive operation.

\begin{enumerate}
    \item Write 0 to the \hyperref[regdesc:AESDMAENABLEREG]{AES\_DMA\_ENABLE\_REG} register before starting the first operation.
    \item Initialize registers \hyperref[regdesc:AESMODEREG]{AES\_MODE\_REG} and \hyperref[regdesc:AESKEYnREG]{AES\_KEY\_\regindex{n}\_REG}
    before starting the first operation.
    \item Update the content of \hyperref[regdesc:AESTEXTINnREG]{AES\_TEXT\_IN\_\regindex{m}\_REG}. \label{others:aes-update-register-typical-continuous}
    \item Start operation by writing 1 to the \hyperref[regdesc:AESTRIGGERREG]{AES\_TRIGGER\_REG} register.
    \item Wait till the content of the \hyperref[regdesc:AESSTATEREG]{AES\_STATE\_REG} register becomes 0, which indicates the operation completes.
    \item Read results from the \hyperref[regdesc:AESTEXTOUTnREG]{AES\_TEXT\_OUT\_\regindex{m}\_REG} register, \label{read result in typical AES continuous calculation} and return to Step \ref{others:aes-update-register-typical-continuous} to continue the next operation.
\end{enumerate}


\section{DMA-AES Working Mode}\label{sec:aes-dma-aes}

In the DMA-AES working mode, the AES accelerator supports six block cipher modes including ECB/CBC/OFB/CTR/CFB8/CFB128. Users can choose the block cipher mode by configuring the \hyperref[regdesc:AESBLOCKMODEREG]{AES\_BLOCK\_MODE\_REG} register according to Table \ref{tab:aes-block-mode} below.

\begin{longtable}[c]{ | c | l | }
\caption{Block Cipher Mode}
\label{tab:aes-block-mode}
\\\hline
\rowcolor{lightgray}
\hyperref[regdesc:AESBLOCKMODEREG]{AES\_BLOCK\_MODE\_REG}[2:0]     & Block Cipher Mode\\
    \hline
    0                          & ECB (Electronic Codebook)    \\
    \hline
    1                          & CBC (Cipher Block Chaining)     \\
    \hline
    2                          & OFB (Output Feedback)    \\
    \hline
    3                          & CTR (Counter)    \\
    \hline
    4                          & CFB8 (8-bit Cipher Feedback)    \\
    \hline
    5                          & CFB128 (128-bit Cipher Feedback) \\
    \hline
    6                         & reserved \\
    \hline
    7                         & reserved \\
    \hline
\end{longtable}

Users can check the working status of the AES accelerator by inquiring the \hyperref[regdesc:AESSTATEREG]{AES\_STATE\_REG} register and comparing the return value against the Table \ref{tab:aes-state-value-dma} below.

\clearpage
\begin{longtable}[c]{ | c | l | l | }
\caption{Working Status under DMA-AES Working mode}
\label{tab:aes-state-value-dma}
\\\hline
\rowcolor{lightgray}
\hyperref[regdesc:AESSTATEREG]{AES\_STATE\_REG}[1:0] & Status &  Description \\ \hline

    0                   & IDLE      &   The AES accelerator is idle.\\ \hline
    1                   & WORK      &   The AES accelerator is in the middle of an operation.     \\ \hline
    2                   & DONE      &   The AES accelerator completed operations.     \\ \hline
\end{longtable}

When working in the DMA-AES working mode, the AES accelerator supports interrupt on the completion of computation. To enable this function, write 1 to the \hyperref[regdesc:AESINTENAREG]{AES\_INT\_ENA\_REG} register. By default, the interrupt function is disabled. Also, note that the interrupt should be cleared by software after use.


\subsection{Key, Plaintext, and Ciphertext}

\textbf{Block Operation}

During the block operations, the AES Accelerator reads source data from DMA, and write result data to DMA after the computation.

\begin{itemize}
    \item For encryption, DMA reads plaintext from memory, then passes it to AES as source data. After computation, AES passes ciphertext as result data  back to DMA to write into memory.
    \item For decryption, DMA reads ciphertext from memory, then passes it to AES as source data. After computation, AES passes plaintext as result data back to DMA to write into memory.
\end{itemize}

During block operations, the lengths of the source data and result data are the same. The total computation time is reduced because the DMA data operation and AES computation can happen concurrently.

\label{api:testpadding}The length of source data for AES Accelerator under DMA-AES working mode must be 128 bits or the integral multiples of 128 bits. Otherwise, trailing zeros will be added to the original source data, so the length of source data equals to the nearest integral multiples of 128 bits. Please see details in Table \ref{tab:aes-define-func-text-padding} below.


\begin{longtable}[c]{|p{15cm}|}
\caption{TEXT-PADDING} 
\label{tab:aes-define-func-text-padding}\\
\hline
\rowcolor{lightgray}
\multirow{1}{*}{%
{\bfseries Function :} {\bfseries TEXT-PADDING}(\ ) } \\\hline
{\bfseries Input\phantom{aaa} :} $X$, bit string.\\
{\bfseries Output\phantom{a~} :} $Y$ = {\bfseries TEXT-PADDING}($X$), whose length is the nearest integral multiples of 128 bits. \\\hline
{\bfseries Steps\phantom{aa~}  }\\
\hspace{1.3cm} Let us assume that X is a data-stream that can be split into {\it{n}} parts as following: \\
\phantom{aaaaaaaa}$X = X_{1} || X_{2} || \cdots || X_{n-1} || X_{n}$ \\
\phantom{aaaaaaaa}Here, the lengths of $X_{1}, X_{2}, \cdots, X_{n-1}$ all equal to 128 bits, and the length of $X_{n}$ is {\it{t}} (0<={\it{t}}<=127).\\
\phantom{aaaaaaaa}If ${\it{t}} = 0$, then\\
\phantom{aaaaaaaa}\phantom{aaaaaaaa} {\bfseries TEXT-PADDING}($X$) = $X$;\\
\phantom{aaaaaaaa}If $0<{\it{t}} <= 127$, define a 128-bit block, $X_{n}^{*}$, and let $X_{n}^{*} = X_{n} || 0^{128-{\it{t}}}$, then\\
\phantom{aaaaaaaa}\phantom{aaaaaaaa} {\bfseries TEXT-PADDING}($X$) = $X_{1} || X_{2} || \cdots || X_{n-1} || X_{n}^{*}$  = $X || 0^{128-{\it{t}}}$\\
\hline
\end{longtable}



\subsection{Endianness}

Under the DMA-AES working mode, the transmission of source data and result data for AES Accelerator is solely controlled by DMA. Therefore, the AES Accelerator cannot control the Endianness of the source data and result data, but does have requirement on how these data should be stored in memory and on the length of the data.

For example, let us assume DMA needs to write the following data into memory at address 0x{}0280.

\begin{itemize}
    \item Data represented in hexadecimal:
    \begin{itemize}
    \item 0102030405060708090A0B0C0D0E0F101112131415161718191A1B1C1D1E1F20
    \end{itemize}
\item Data Length:
    \begin{itemize}
    \item Equals to 2 blocks.
    \end{itemize}
\end{itemize}

Then, this data will be stored in memory as shown in Table \ref{tab:aes-dma-aes-plaintext-ciphertext-endian} below.

\begin{longtable}[c]
{ | l | l |l | l | l | l | l | l |}
\caption{Text Endianness for DMA-AES} 
\label{tab:aes-dma-aes-plaintext-ciphertext-endian}\\
\hline
\rowcolor{lightgray}
Address &  Byte  & Address &  Byte  & Address &  Byte  & Address &  Byte \\\hline
\endfirsthead
\hline
\rowcolor{lightgray}
Address &  Byte  & Address &  Byte  & Address &  Byte  & Address &  Byte \\\hline
\endhead
\endfoot
0x{}0280 & 0x{}01 & 0x{}0281 & 0x{}02 & 0x{}0282 & 0x{}03 & 0x{}0283 & 0x{}04 \\\hline
0x{}0284 & 0x{}05 & 0x{}0285 & 0x{}06 & 0x{}0286 & 0x{}07 & 0x{}0287 & 0x{}08 \\\hline
0x{}0288 & 0x{}09 & 0x{}0289 & 0x{}0A & 0x{}028A & 0x{}0B & 0x{}028B & 0x{}0C \\\hline
0x{}028C & 0x{}0D & 0x{}028D & 0x{}0E & 0x{}028E & 0x{}0F & 0x{}028F & 0x{}10 \\\hline
0x{}0290 & 0x{}11 & 0x{}0291 & 0x{}12 & 0x{}0292 & 0x{}13 & 0x{}0293 & 0x{}14 \\\hline
0x{}0294 & 0x{}15 & 0x{}0295 & 0x{}16 & 0x{}0296 & 0x{}17 & 0x{}0297 & 0x{}18 \\\hline
0x{}0298 & 0x{}19 & 0x{}0299 & 0x{}1A & 0x{}029A & 0x{}1B & 0x{}029B & 0x{}1C \\\hline
0x{}029C & 0x{}1D & 0x{}029D & 0x{}1E & 0x{}029E & 0x{}1F & 0x{}029F & 0x{}20 \\\hline
\end{longtable}


DMA can access both internal memory and PSRAM outside \chipname{}. When you use DMA to access external PSRAM, please use base addresses that meet the requirements for DMA. When you use DMA to access internal memory,  base addresses do not have such requirements. Details can be found in Chapter \ref{mod:dma} \textit{\nameref{mod:dma}}.


\subsection{Standard Incrementing Function}\label{sec:aes-Standard-Incrementing-Function}

AES accelerator provides two Standard Incrementing Functions for the CTR block operation, which are INC$_{32}$ and INC$_{128}$ Standard Incrementing Functions. By setting the  \hyperref[regdesc:AESINCSELREG]{AES\_INC\_SEL\_REG} register to 0 or 1, users can choose the INC$_{32}$ or INC$_{128}$ functions respectively. For details on the Standard Incrementing Function, please see Chapter B.1 The Standard Incrementing Function in \href{https://nvlpubs.nist.gov/nistpubs/Legacy/SP/nistspecialpublication800-38a.pdf}{NIST SP 800-38A}.

\subsection{ Block Number}\label{sec:aes-block-number}

Register \hyperref[regdesc:AESBLOCKNUMREG]{AES\_BLOCK\_NUM\_REG} stores the Block Number of plaintext $P$ or ciphertext $C$. The length of this register equals to length({\bfseries TEXT-PADDING}($P$))/128 or length({\bfseries TEXT-PADDING}($C$))/128. The AES Accelerator only uses this register when working in the DMA-AES mode.

\subsection{Initialization Vector}\label{sec:aes-init-vector}

\label{memdesc:AESIVMEM}\hyperref[memdesc:AESIVMEM]{AES\_IV\_MEM} is a 16-byte memory, which is only available for AES Accelerator working in block operations. For CBC/OFB/CFB8/CFB128 operations, the \hyperref[memdesc:AESIVMEM]{AES\_IV\_MEM} memory stores the Initialization Vector (IV). For the CTR operation, the \hyperref[memdesc:AESIVMEM]{AES\_IV\_MEM} memory stores the Initial Counter Block (ICB).

Both IV and ICB are 128-bit strings, which can be divided into Byte0, Byte1, Byte2 \begin{math}\cdots\end{math} Byte15 (from left to right). \hyperref[memdesc:AESIVMEM]{AES\_IV\_MEM} stores data following the Endianness pattern presented in Table \ref{tab:aes-dma-aes-plaintext-ciphertext-endian}, i.e., the most significant (i.e., left-most) byte Byte0 is stored at the lowest address while the least significant (i.e., right-most) byte Byte15 at the highest address.

For more details on IV and ICB, please refer to \href{https://nvlpubs.nist.gov/nistpubs/Legacy/SP/nistspecialpublication800-38a.pdf}{NIST SP 800-38A}.

\subsection{ Block Operation Process}
\begin{enumerate}
    \item Select one of DMA channels to connect with AES, configure the DMA chained list, and then start DMA. For details, please refer to Chapter \ref{mod:dma} \textit{\nameref{mod:dma}}.
    \item Initialize the AES accelerator-related registers:
    \begin{itemize}
        \item Write 1 to the \hyperref[regdesc:AESDMAENABLEREG]{AES\_DMA\_ENABLE\_REG} register.
        \item Configure the \hyperref[regdesc:AESINTENAREG]{AES\_INT\_ENA\_REG} register to enable or disable the interrupt function.
        \item Initialize registers \hyperref[regdesc:AESMODEREG]{AES\_MODE\_REG} and \hyperref[regdesc:AESKEYnREG]{AES\_KEY\_\regindex{n}\_REG}.
        \item Select block cipher mode by configuring the \hyperref[regdesc:AESBLOCKMODEREG]{AES\_BLOCK\_MODE\_REG} register. For details, see Table \ref{tab:aes-block-mode}.
        \item Initialize the \hyperref[regdesc:AESBLOCKNUMREG]{AES\_BLOCK\_NUM\_REG} register. For details, see Section \ref{sec:aes-block-number}.
        \item Initialize the \hyperref[regdesc:AESINCSELREG]{AES\_INC\_SEL\_REG} register (only needed when AES Accelerator is working under CTR block operation).
        \item Initialize the \hyperref[memdesc:AESIVMEM]{AES\_IV\_MEM} memory (This is always needed except for ECB block operation).
    \end{itemize}
    \item Start operation by writing 1 to the \hyperref[regdesc:AESTRIGGERREG]{AES\_TRIGGER\_REG} register.
    \item \label{other:aes-wait-ECB-CBC-OFB-CTR-CFB8-CFB128-done}
            Wait for the completion of computation, which happens when the content of \hyperref[regdesc:AESSTATEREG]{AES\_STATE\_REG} becomes 2 or the AES interrupt occurs.
    \item Check if DMA completes data transmission from AES to memory. At this time, DMA had already written the result data in memory, which can be accessed directly. For details on DMA, please refer to Chapter \ref{mod:dma} \textit{\nameref{mod:dma}}.
    \item Clear interrupt by writing 1 to the \hyperref[regdesc:AESINTCLRREG]{AES\_INT\_CLR\_REG} register, if any AES interrupt occurred during the computation.
    \item Release the AES Accelerator by writing 0 to the \hyperref[regdesc:AESDMAEXITREG]{AES\_DMA\_EXIT\_REG} register. After this, the content of the \hyperref[regdesc:AESSTATEREG]{AES\_STATE\_REG} register becomes 0.
    Note that, you can release DMA earlier, but only after Step \ref{other:aes-wait-ECB-CBC-OFB-CTR-CFB8-CFB128-done} is completed.
\end{enumerate}

\section{Memory Summary}\label{aesMemorySummary}

The addresses in this section are relative to the AES accelerator base address provided in Table \ref{tab:sysmem-base-address} in Chapter \ref{mod:sysmem} \textit{\nameref{mod:sysmem}}.

The abbreviations given in Column \textbf{Access} are explained in Section \hyperref[glossary-access-types]{\textit{Access Types for Registers}}.

\subfile{\modulefiles/30-AES-mem-summ__EN.tex}


\clearpage
\section{Register Summary} \hypertarget{aes-reg-summ}{}
The addresses in this section are relative to the AES accelerator base address provided in Table \ref{tab:sysmem-base-address} in Chapter \ref{mod:sysmem} \textit{\nameref{mod:sysmem}}.

The abbreviations given in Column \textbf{Access} are explained in Section \hyperref[glossary-access-types]{\textit{Access Types for Registers}}.

\subfile{\modulefiles/30-AES-reg-summ__EN.tex}

\clearpage
\section{Registers}
The addresses in this section are relative to the AES accelerator base address provided in Table \ref{tab:sysmem-base-address} in Chapter \ref{mod:sysmem} \textit{\nameref{mod:sysmem}}.

\subfile{\modulefiles/30-AES-reg__EN.tex}

\end{document}
